% git_hash=ce666856, timestamp=2026-07-16 16:34:28 EDT
%%%% Chapter file for Integrating Causality and Probability in ML %%%%
% This chapter file can be compiled standalone or included in the root book.tex

\chapter{Integrating Causality and Probability in ML}
\label{chap:integrating-causality-and-probability-in-ml}

\motto{A prediction says what will happen. A decision says what to do.}

% From: book_springer/lectures_source/Lesson02.2_Integrating_Causality_And_Probability_in_ML.smd:17 '# Bridging Prediction and Decision'
\section{Bridging Prediction and Decision}
\label{sec:bridging-prediction-and-decision}

% From: book_springer/lectures_source/Lesson02.2_Integrating_Causality_And_Probability_in_ML.smd:19 '## From Costs to Tools'
\subsection{From Costs to Tools}
\label{subsec:from-costs-to-tools}

The previous chapter left four costs unaddressed: the cost of ignoring causality,
the cost of ignoring uncertainty, the cost of ignoring the business objective, and
the cost of ignoring dynamics. This lesson builds the machinery needed to
address the first two, starting with the vocabulary that makes a causal claim precise
enough to act on. The business framing established earlier does not change: \textit{a
prediction says what will happen, and a decision says what to do}. What this chapter
adds is the machinery sitting underneath that framing, the apparatus that turns
a causal claim from a slogan into a quantity that can be computed from data

% From: book_springer/lectures_source/Lesson02.2_Integrating_Causality_And_Probability_in_ML.smd:19 '* From Costs to Tools: What This Chapter Builds'
\textbf{From Costs to Tools: What This Chapter Builds} --- The previous chapter
catalogued four costs of ignoring causality, uncertainty, the business objective,
and dynamics in a decision system. This lesson does not revisit those failures
in detail; instead, it builds the tools that close the first two gaps. The business
framing stays fixed throughout: a prediction says what will happen, and a decision
says what to do. What this chapter adds is the machinery that makes causal
claims precise enough to act on

This chapter's contribution organizes around three pieces of apparatus. The
first is \emph{the ladder of causation}, which formalizes what the informal
phrases ``what if'' and ``why'' actually mean, separating association,
intervention, and counterfactual reasoning into distinct rungs with distinct evidence
requirements. The second is \emph{causal graphs and identification}, which turn
a causal question stated in words—``does price cuts boost sales?''—into a computable
quantity, an expression built from the observable variables a dataset actually
contains. The third is \emph{a probabilistic toolkit}, which turns a single point
estimate into a posterior distribution a decision-maker can trust, replacing a
bare number with a quantified sense of how confident that number deserves to be

Each piece answers a different half of the causality-and-uncertainty gap: the
ladder and the graphs supply the causal half, saying what a claim like ``if we cut
price 10\%, revenue rises'' would even mean and what evidence would settle it;
the probabilistic toolkit supplies the uncertainty half, attaching a defensible range
to whatever the causal machinery concludes. Neither half is optional. A causal estimate
reported without uncertainty is the point-estimate failure mode from the previous
chapter, dressed up in causal notation; an uncertainty estimate computed over the
wrong, purely observational quantity is the correlation-versus-causation failure
mode, dressed up in confidence intervals. The rest of this lesson works through
each piece of the apparatus in turn, beginning with the ladder of causation

% From: book_springer/lectures_source/Lesson02.2_Integrating_Causality_And_Probability_in_ML.smd:36 '* Prediction vs Decision, Formally'
\textbf{Prediction vs Decision, Formally} --- Prediction and decision are not the
same kind of mathematical object, and treating them as interchangeable is a
common source of costly errors, exactly the kind of error the ice-cream-and-drowning
and loan-approval cases explored in the previous chapter. Table~\ref{tab:prediction-vs-decision-formally}
sets the two side by side across five dimensions

\begin{table}[!t]
  \caption{Prediction and decision compared across target, data, metric, assumption,
  and how each treats the world.}
  \label{tab:prediction-vs-decision-formally}
  \begin{tabular}{p{2.0cm}p{4.5cm}p{4.8cm}}
    \hline
    \noalign{\smallskip} Aspect                             & Prediction             & Decision                                  \\
    \noalign{\smallskip}\svhline\noalign{\smallskip} Target & $\mathbb{E}[Y \mid X]$ & $\arg\max_{a}\mathbb{E}[U(Y) \mid do(a)]$ \\
    Data                                                    & Observational only     & Observational + experimental              \\
    Metric                                                  & Accuracy, AUC, MSE     & Regret, expected utility                  \\
    Assumption                                              & i.i.d., stable world   & Models the intervention                   \\
    World                                                   & No feedback            & Actions change the world                  \\
    \noalign{\smallskip}                                     \hline
    \noalign{\smallskip}
  \end{tabular}
\end{table}

Reading down the columns of Table~\ref{tab:prediction-vs-decision-formally}, a predictive
target is a conditional expectation, $\mathbb{E}[Y \mid X]$, estimated from whatever
data happened to be logged, evaluated by how closely it tracks the outcome that
actually occurred, and reported under the assumption that the data-generating process
holds still. A decision target is an optimization over actions, $\arg\max_{a}\mathbb{E}
[U(Y) \mid do(a)]$, and every one of those properties changes with it: it needs data
on what happens \emph{under} an action, not merely data on what an action happened
to correlate with; it is scored by regret or expected utility rather than by
accuracy; and it explicitly models the intervention rather than assuming the world
stays fixed while it is measured. The bottom row of the table is the same
warning explored at length under dynamics and feedback: a decision, once acted
on, changes the very world the next prediction will be scored against, while a
prediction, by construction, assumes it will not

The key insight is that a decision target contains a $do(\cdot)$ that a predictive
target never does, and essentially everything introduced in this lesson—the
ladder of causation, causal graphs, identification, the probabilistic toolkit—exists
for one purpose: computing that $do(\cdot)$ from data that was, in the overwhelming
majority of real settings, never generated by an experiment in the first place

% From: book_springer/lectures_source/Lesson02.2_Integrating_Causality_And_Probability_in_ML.smd:58 '# The Ladder of Causation'
\section{The Ladder of Causation}
\label{sec:the-ladder-of-causation}

% From: book_springer/lectures_source/Lesson02.2_Integrating_Causality_And_Probability_in_ML.smd:60 '## Three Rungs of Causal Reasoning'
\subsection{Three Rungs of Causal Reasoning}
\label{subsec:three-rungs-of-causal-reasoning}

The apparatus promised in the previous section starts with a map of the different
kinds of causal question a model might be asked to answer, and of what each kind
of question actually demands as evidence. This map is Pearl's ladder of causation,
which organizes causal reasoning into three rungs of increasing power

% From: book_springer/lectures_source/Lesson02.2_Integrating_Causality_And_Probability_in_ML.smd:62 '* The Ladder of Causation'
\textbf{The Ladder of Causation} --- Judea Pearl organizes causal reasoning into
three rungs of increasing power, each one strictly more demanding than the last in
terms of what it asks a model to know. Table~\ref{tab:ladder-of-causation} lays
out the three rungs side by side, with the symbol, the everyday activity, and the
question each rung answers

\begin{table}[!t]
  \caption{Pearl's ladder of causation: three rungs of increasing power, from
  seeing to doing to imagining.}
  \label{tab:ladder-of-causation}
  \begin{tabular}{p{2.3cm}p{3.4cm}p{2.0cm}p{2.6cm}}
    \hline
    \noalign{\smallskip} Rung                                       & Symbol                  & Activity  & Question \\
    \noalign{\smallskip}\svhline\noalign{\smallskip} 1. Association & $\Pr(Y \mid X)$         & Seeing    & What is? \\
    2. Intervention                                                 & $\Pr(Y \mid do(X), Z)$  & Doing     & What if? \\
    3. Counterfactual                                               & $\Pr(Y_{x}\mid x', y')$ & Imagining & Why?     \\
    \noalign{\smallskip}                                             \hline
    \noalign{\smallskip}
  \end{tabular}
\end{table}

The rungs in Table~\ref{tab:ladder-of-causation} are not three equally valid ways
of describing the same reasoning; they are ordered, and climbing from one to the
next requires strictly more than a bigger dataset of the same kind. Rung~1, association,
asks what seeing $X$ reveals about $Y$, and it is answered entirely from observational
data. Rung~2, intervention, asks what happens if $X$ is set by an outside action,
$do(X)$, rather than merely observed, holding a set of other variables $Z$ fixed,
and answering it generally requires either an experiment or a causal model built
from assumptions that no amount of passive observation can verify on its own. Rung~3,
counterfactual reasoning, asks the sharpest question of the three: what would
have happened to a specific unit, characterized by an observed treatment $x'$ and
outcome $y'$, had a different treatment $x$ been assigned instead This is a
strictly harder question than the population-level intervention question, since it
requires reasoning about a world that, for that particular unit, never occurred
The critical point is that a model fitting only $\Pr(Y \mid X)$ lives on rung~1,
and no amount of additional rung-1 data, however large the dataset, answers a rung-2
or rung-3 question, because the rungs differ in kind, not in degree; more
association-level evidence never becomes intervention-level or counterfactual-level
evidence on its own

% From: book_springer/lectures_source/Lesson02.2_Integrating_Causality_And_Probability_in_ML.smd:81 '* Rung 1: Association'
\textbf{Rung 1: Association} --- Rung~1 asks a single, intuitive question: how
should seeing $X$ change a belief about $Y$? Formally, this is the conditional
distribution $\Pr(Y \mid X)$, and answering it requires nothing more than passive
observation of whether $X$ and $Y$ tend to occur together. This is also the
entire territory traditional predictive machine learning occupies: a classifier,
a regressor, or a forecasting model, however sophisticated its architecture, is fitting
some version of $\Pr(Y \mid X)$ and stopping there

Two everyday examples make the rung concrete. A symptom is informative about a disease
in exactly the rung-1 sense: observing the symptom shifts a belief about which
disease is present, without anyone needing to know or model the biological mechanism
connecting the two. A poll is informative about an election result the same way,
since seeing the poll numbers shifts a belief about the likely outcome purely because
the two have tended to move together historically, not because the poll causes
the result

The rung has a sharp limitation, and it is the same limitation the previous
chapter spent an entire section on: association alone cannot distinguish a true cause
from a confounder, and it cannot distinguish a true cause from reverse causation
Seeing that a symptom and a disease move together does not, by itself, say
whether the symptom causes complications, whether the disease causes the symptom,
or whether some third factor produces both. A model that lives entirely on rung~1
has no way to adjudicate between these explanations, because all of them predict
the identical observational pattern; rung-1 evidence is, in this precise sense, silent
on the question that rungs~2 and~3 exist to answer. This is precisely the gap
identified from the business side—confounding and selection bias masquerading as
causal effects—and it is why the next rungs of the ladder, not a larger rung-1 dataset,
are the tools that close it

% From: book_springer/lectures_source/Lesson02.2_Integrating_Causality_And_Probability_in_ML.smd:93 '* Rung 2: Intervention'
\textbf{Rung 2: Intervention} --- Rung~2 asks a sharper question: what happens to
$Y$ if I do $X$? The mathematical notation $\Pr(Y \mid do(X), Z)$ looks similar
to $\Pr(Y \mid X, Z)$, but the $do(\cdot)$ operator marks a profound difference Setting
$X$ by external action, not by observation, requires a causal model or a controlled
experiment, because the question is fundamentally about what the world would
look like if an action were taken, not merely what the data happens to show when
$X$ and $Y$ are observed together

Two concrete examples illustrate the difference from rung~1. If we observe that
stores offering a discount sell more units, that is rung-1 reasoning: the
discount and sales are associated. But the interventional question—if we cut price
10\%, what happens to units sold?—requires knowing whether the correlation between
discount and sales reflects a causal effect or whether, for instance, discounts are
offered only when inventory is high anyway. Similarly, if we ban sodas from schools,
what happens to obesity? The answer could differ sharply from what observation of
soda consumption and obesity patterns suggests, because the causal structure
underlying the relationship is what matters, not the correlation

\textbf{Key idea}: $do(X)$ severs $X$ from its usual causes, so the answer can differ
sharply from what observation suggests. This difference is not a matter of
having more data; it is a matter of the kind of data and the assumptions layered
on top of it. A randomized trial gives rung-2 answers for free, by ensuring that
$X$ depends on nothing but the random assignment. Observational data can give
rung-2 answers too, but only if the causal structure is known and certain conditions—the
identification assumptions introduced later in this chapter—are met

% From: book_springer/lectures_source/Lesson02.2_Integrating_Causality_And_Probability_in_ML.smd:105 '* Rung 3: Counterfactuals'
\textbf{Rung 3: Counterfactuals} --- Rung~3 is the highest, sharpest, and most demanding
rung of the ladder. It asks: was it $X$ that caused $Y$, for this specific unit?
The mathematical notation $\Pr(Y_{x}\mid x', y')$ expresses the outcome had $X$ been
set to $x$, given that we observed $X = x'$ and outcome $Y = y'$ for this
particular unit. Counterfactual reasoning is about imagining an alternative to what
actually happened; for a given individual, exactly one of the two potential outcomes
(treatment or control) is observed, and the other is forever counterfactual

The conceptual gap between rung~2 and rung~3 is the gap between a population-level
answer and a unit-level answer. Rung~2 asks: does this intervention work on average?
Rung~3 asks: would this intervention have worked for this person? A business can
observe that a customer who opened a promotional email bought more, but rung~3 reasoning
answers the harder question: would this customer have bought at all if we had
not sent the email? This is the language of attribution, blame, and credit

Examples of rung~3 questions abound in practice. Would this customer have churned
had we offered a discount? Did the training program cause this employee's
promotion, or would the promotion have happened anyway? Why did the campaign
fail? Each of these questions requires a full structural model that can simulate
not just interventions, but specific counterfactual histories for specific units
Rung~3 is also the rung most demanding of data and assumptions, since answering it
requires both identification (like rung~2) and an estimate of the conditional distribution
of the unobserved potential outcome, given what was observed

% From: book_springer/lectures_source/Lesson02.2_Integrating_Causality_And_Probability_in_ML.smd:117 '* Seeing vs Doing: A Worked Example'
\textbf{Seeing vs Doing: A Worked Example} --- To make the ladder concrete, consider
a tornado-warning system. A tornado forms (denoted $T$). An official warning is issued
($W$). Radio and TV broadcast the warning ($A$ and $B$) Because both channels
wait for the official order, the warning causes both broadcasts: $W \to A$ and
$W \to B$. Residents are warned ($R$) if either channel broadcasts:
$R = A \lor B$

On rung~1 (seeing), we observe the data: whenever radio $A$ broadcasts, TV $B$
also broadcasts. The association is real and detectable from data. An observer
seeing $A = 1$ would update their belief that $B = 1$ is higher, because they co-occur
through the common cause $W$

On rung~2 (doing), we perform $do(A = 1)$: we force the radio to broadcast. This
severs the edge $W \to A$. Now the mechanism is different. Radio broadcast no longer
depends on whether the official warning was issued, because we have forced it
externally. With $A$ forced to 1, TV broadcast $B$ is unaffected, because the only
link between them was through $W$, which we have now disconnected from $A$. So $d
o(A=1)$ makes TV broadcast less likely—because if radio is forced on, the
official order is less likely, so TV is less likely to receive the broadcast order
The conclusion is opposite to what rung~1 suggested

This example shows why seeing and doing can give opposite answers. Seeing $A$ and
doing $A$ are not the same because observation updates belief about the causes of
$A$, while intervention severs the causes of $A$

% Figure: Tornado warning DAG
% rendered_image:begin
% filename: figures/tornado_warning_dag.png
% label: fig:tornado-warning
% caption: Tornado warning system: $T$ (tornado forms) causes $W$ (warning issued), which causes both $A$ (radio broadcast) and $B$ (TV broadcast). Both $A$ and $B$ cause $R$ (residents warned). Seeing $A$ and $B$ co-occur via $W$; doing $do(A=1)$ severs the $W \to A$ edge.
% rendered_image:end

% From: book_springer/lectures_source/Lesson02.2_Integrating_Causality_And_Probability_in_ML.smd:162 '# Causal Graphical Models'
\section{Causal Graphical Models}
\label{sec:causal-graphical-models}

% From: book_springer/lectures_source/Lesson02.2_Integrating_Causality_And_Probability_in_ML.smd:164 '## Structure and Variable Types'
\subsection{Structure and Variable Types}
\label{subsec:structure-and-variable-types}

% From: book_springer/lectures_source/Lesson02.2_Integrating_Causality_And_Probability_in_ML.smd:166 '* Structural Causal Models'
\textbf{Structural Causal Models} --- The ladder of causation is the conceptual framework;
causal graphs are the machinery. A \emph{structural causal model} (SCM) writes
each variable as a function of its direct causes and an exogenous noise term:
\begin{equation}
  X_{i}= f_{i}(\text{Parents}(X_{i}), \varepsilon_{i}) \label{eq:scm-definition}
\end{equation}

This is more than a notational trick. An SCM answers all three rungs from one object:
observation ($\Pr(Y \mid X = x)$), intervention ($\Pr(Y \mid do(X = x))$), and
counterfactual ($\Pr(Y_{X=x}\mid X = x')$). The key is that the structural equations,
not just the shape of the graph, are what let an SCM simulate an intervention
When we $do(X = x)$, we replace the structural equation $X = f_{X}(\text{Parents}
(X), \varepsilon_{X})$ with the fixed value $x$. The graph tells us which edges
to delete (all incoming edges to $X$); the structural equations tell us how the rest
of the system evolves

A concrete example makes this clear. Consider a garden network: whether it is
cloudy ($C$) drives both whether the sprinkler activates ($S$) and whether it
rains ($R$). Rain and sprinkler use both drive whether the grass gets wet ($W$) The
factorization is $\Pr(C, R, S, W) = \Pr(W \mid R, S) \Pr(S \mid C) \Pr(R \mid C)
\Pr(C)$. If we observe cloudy weather and see the grass is wet, we update our belief
about whether it rained or the sprinkler was on. But if we do $(S = 1)$ (force the
sprinkler on), we do not update our belief about cloudiness—the sprinkler does not
change the weather. The structural equations let the SCM make this distinction
automatically Crucially, the structural equations, not just the shape of the
graph, are what let an SCM simulate an intervention. The graph structure alone
is not enough; the functional form and independence assumptions embedded in Equation~\eqref{eq:scm-definition}
are what make intervention, not just observation, computable

% Figure: Sprinkler DAG
% rendered_image:begin
% filename: figures/sprinkler_dag.png
% label: fig:sprinkler-dag
% caption: Garden network: $C$ (cloudy) drives $S$ (sprinkler) and $R$ (rain); both drive $W$ (wet grass).
% rendered_image:end

% From: book_springer/lectures_source/Lesson02.2_Integrating_Causality_And_Probability_in_ML.smd:202 '* Confounders: The Fork Structure'
\textbf{Confounders: The Fork Structure} --- One of the most important causal structures
is the fork: $X \leftarrow Z \rightarrow Y$, where $Z$ is a \emph{confounder}, a
common cause of both $X$ and $Y$. Confounding is what happens when $X$ and $Y$
are dependent even though neither causes the other; they are dependent because they
share a common cause

A concrete example: does hiring expensive consultants improve future profits?
Observationally, firms that hire consultants do report higher future profits than
those that do not. But here is the confounder: past profits drive the decision to
hire consultants (firms with spare cash invest in consultants), and past profits
also drive future profits (firms with momentum tend to stay ahead). So past profits
is the common cause, and the observed correlation between consultant hiring and
future profits is spurious—it is confounding, not causation

The key insight is that conditioning on $Z$ blocks the spurious association. If we
compare consultant-hiring firms to non-hiring firms that had identical past profits,
the spurious effect disappears. Formally, $(X \perp Y) \mid Z$ (conditioned on $Z$,
$X$ and $Y$ become independent). This is the insight behind back-door adjustment,
developed later: a predictive model that conditions on the confounder will not
bake in the spurious association. Without conditioning on $Z$, however, a predictive
model bakes in the spurious association between $X$ and $Y$—a form of selection
bias or confounding that can masquerade as a causal effect and is one of the
four costs of ignoring causality that the previous chapter identified

% Figure: Confounder DAG
% rendered_image:begin
% filename: figures/confounder_dag.png
% label: fig:confounder
% caption: Fork structure: Past profits ($Z$) confound the relationship between hiring consultants ($X$) and future profits ($Y$).
% rendered_image:end

% From: book_springer/lectures_source/Lesson02.2_Integrating_Causality_And_Probability_in_ML.smd:236 '* Mediators: The Chain Structure'
\textbf{Mediators: The Chain Structure} --- Another fundamental structure is the
chain: $X \rightarrow M \rightarrow Y$, where $M$ is a \emph{mediator} that
transmits the causal effect of $X$ on $Y$. Unlike a confounder, a mediator is
\emph{on} the causal path from $X$ to $Y$; it is the mechanism through which $X$
affects $Y$

Suppose a training program ($X$) improves productivity ($Y$) through job satisfaction
($M$). The training raises satisfaction, and higher satisfaction raises productivity
There are two effects: the \emph{indirect effect} of training through
satisfaction, and possibly a \emph{direct effect}, any remaining path from training
to productivity not through satisfaction (e.g., the training teaches new skills
directly). The total effect of training on productivity is the sum of direct and
indirect paths

Crucially, the treatment of a mediator differs from that of a confounder. When estimating
the total causal effect of $X$ on $Y$, a mediator should \emph{not} be conditioned
on, because conditioning on $M$ blocks the indirect effect, making the total effect
appear smaller than it is. If we condition on job satisfaction when estimating
the effect of training on productivity, we remove the path through satisfaction,
leaving only the direct path and underestimating the total causal benefit. Unlike
a confounder, a mediator should not be conditioned on when the goal is the total
effect of $X$ on $Y$; conditioning on a mediator can introduce bias just as surely
as omitting a confounder, but in the opposite direction

% Figure: Mediator DAG
% rendered_image:begin
% filename: figures/mediator_dag.png
% label: fig:mediator
% caption: Chain structure: Training program ($X$) affects productivity ($Y$) through job satisfaction ($M$). The dashed line represents a possible direct effect.
% rendered_image:end

% From: book_springer/lectures_source/Lesson02.2_Integrating_Causality_And_Probability_in_ML.smd:271 '* Colliders Revisited: A Formal Treatment'
\textbf{Colliders Revisited: A Formal Treatment} --- A \emph{collider} is the
structure $X \rightarrow M \leftarrow Y$, where $M$ is influenced by both $X$ and
$Y$. At first glance, colliders seem like they should open a path between $X$ and
$Y$: both arrows point into $M$. But the opposite is true: $X$ and $Y$ are
initially independent when $M$ is not conditioned on. Conditioning on $M$, or on
any descendant of $M$, makes $X$ and $Y$ dependent, a phenomenon called \emph{explaining
away}

Consider exercise ($X$) and diet ($D$) as two factors affecting body weight ($W$),
which in turn affects heart disease ($H$). Exercise and diet are independent in
the general population; people who exercise more are not inherently more likely to
eat healthier. But if we condition on body weight—say, we look only at people at
a given body weight—the independence breaks. Among people at a fixed body weight,
less exercise implies a stricter diet, and vice versa. The reason is that, given
body weight, the model must explain that body weight using either exercise or diet,
so observing one makes the other more likely

Formally, a path is blocked given a conditioning set $Z$ if it contains a non-collider
in $Z$, or a collider with no descendant in $Z$. Two variables are \emph{d-separated}
(and hence independent) if every path between them is blocked. The collider rule
is: a collider blocks a path unless we condition on the collider or one of its descendants
Conditioning on a collider or its descendant creates spurious association—a
different form of bias from confounding, and easy to miss in practice. It looks
like "I just need to condition on more variables," but conditioning on the wrong
variables (descendants of colliders) can create bias instead of removing it

% Figure: Collider DAG
% rendered_image:begin
% filename: figures/collider_dag.png
% label: fig:collider
% caption: Collider structure: Exercise ($E$) and diet ($D$) both affect body weight ($W$), which affects heart disease ($H$). Among people with a fixed body weight, exercise and diet are no longer independent.
% rendered_image:end

% From: book_springer/lectures_source/Lesson02.2_Integrating_Causality_And_Probability_in_ML.smd:310 '* Building a Causal DAG from Domain Knowledge'
\textbf{Building a Causal DAG from Domain Knowledge} --- The structures above (forks,
chains, colliders) are building blocks. The real work of causal inference begins
with translating domain knowledge into a causal graph. Nodes are like nouns (price,
sales, revenue); edges are like verbs (causes, mediates). Every variable must be
classified on two dimensions:

\begin{itemize}
  \item \emph{Observed or unobserved}: Is the variable measured in the data?

  \item \emph{Endogenous or exogenous}: Does the model explain it, or is it
    taken as given?
\end{itemize}

Consider modeling income. Economic policy is exogenous and observed (we know
what policy is in effect, and the model does not explain it). Natural talent is exogenous
and unobserved (the model does not explain it, and we do not measure it directly)
Education and income are endogenous and observed (the model explains both, and
both are measured). The unobserved confounder—natural talent—still has to be
represented in the graph as a dashed node with outgoing edges. If it is omitted entirely,
the identification step later will silently assume it away, and the resulting causal
estimate will be biased

Domain experts' judgments about causality are built into the graph structure Different
experts might draw different graphs for the same domain; the graph is a
hypothesis about the causal system, and no dataset alone can verify it. What the
data can do is verify whether the consequences of the graph (the independencies implied
by d-separation) are consistent with the observed correlations. An unobserved confounder
does not disappear because it is unmeasured; it still has to be represented in
the graph, or the identification step will silently assume it away. The graph is
a quantitative hypothesis about causality, and building it carefully,
incorporating domain knowledge and expert judgment, is as important as the formal
machinery that follows

% From: book_springer/lectures_source/Lesson02.2_Integrating_Causality_And_Probability_in_ML.smd:329 '# Identifying Causal Effects'
\section{Identifying Causal Effects}
\label{sec:identifying-causal-effects}

% From: book_springer/lectures_source/Lesson02.2_Integrating_Causality_And_Probability_in_ML.smd:331 '## From Association to Intervention'
\subsection{From Association to Intervention}
\label{subsec:from-association-to-intervention}

The graph supplies the qualitative structure; now we need to know: given this
structure and observational data, can we compute $\Pr(Y \mid do(X))$? This is
the identification problem, and it has three main answers, each suited to
different graph structures

% From: book_springer/lectures_source/Lesson02.2_Integrating_Causality_And_Probability_in_ML.smd:333 '* The Do-Operator: Simulating Intervention'
\textbf{The Do-Operator: Simulating Intervention} --- The $do(\cdot)$ operator
is how we formally represent an intervention. To compute $\Pr(Y \mid do(X = x))$,
we replace $X$'s structural equation with the fixed value $x$. Formally, this
deletes every incoming edge to $X$ in the causal graph, producing a \emph{mutilated}
graph The key insight is that observing $X = x$ updates our belief about $X$'s
causes, while $do(X = x)$ severs them:
\begin{equation}
  \Pr(Y \mid do(X = x)) \neq \Pr(Y \mid X = x) \label{eq:do-vs-observe}
\end{equation}

In the sprinkler example, $do(\text{Sprinkler}= \text{true})$ removes the edge
$\text{Cloudy}\to \text{Sprinkler}$. Only the descendants of Sprinkler (here, WetGrass)
are affected; Cloudy and Rain are untouched. If we force the sprinkler on, we have
no reason to update our belief about whether it is cloudy, because we have taken
the decision out of cloudy weather's hands. The do-operator turns a graph into a
simulator of interventions the business has not yet tried—the formal machinery behind
counterfactual prediction: what would happen if we did this?

% Figure: Do-operator mutilated graph
% rendered_image:begin
% filename: figures/do_operator_dag.png
% label: fig:do-operator
% caption: Mutilated graph after $do(\text{Sprinkler} = \text{true})$: The edge $\text{Cloudy} \to \text{Sprinkler}$ is removed, severing the sprinkler from weather effects.
% rendered_image:end

% From: book_springer/lectures_source/Lesson02.2_Integrating_Causality_And_Probability_in_ML.smd:369 '* Randomized Controlled Trials: The Gold Standard'
\textbf{Randomized Controlled Trials: The Gold Standard} --- The most
straightforward way to compute $\Pr(Y \mid do(X))$ is with a randomized
controlled trial (RCT). Randomizing $X$ is the physical instantiation of the do-operator:
random assignment severs every incoming edge into $X$, known and unknown alike It
turns the observational $\Pr(Y \mid X)$ into a valid estimate of $\Pr(Y \mid do(X
))$ automatically

Consider assigning students to a tutoring program. If we let students opt in,
self-selection confounds the effect: motivated students both opt in and perform better,
so observing correlation between tutoring and performance does not tell us
whether tutoring caused the improvement. Randomizing the assignment removes this
confounder by severing the link between motivation and assignment. Now comparing
the outcomes of treated and control students gives an unbiased estimate of the
causal effect

Randomized trials have limitations. They may be unethical (randomizing exposure to
a harmful treatment), expensive, slow, or subject to non-compliance and attrition
(people assigned to treatment drop out). When a trial is not possible, the graph
and identification criteria have to supply what a trial would have given for free
This is why identification, the topic of the next sections, is essential: it tells
us when and how we can compute $\Pr(Y \mid do(X))$ from purely observational data

% From: book_springer/lectures_source/Lesson02.2_Integrating_Causality_And_Probability_in_ML.smd:386 '* The Back-Door Criterion'
\textbf{The Back-Door Criterion} --- The fundamental identification question is:
can $\Pr(Y \mid do(X))$ be computed from purely observational data, without a randomized
trial? The back-door criterion answers this by identifying which variables must
be conditioned on

A \emph{back-door path} from $X$ to $Y$ is any path starting with an arrow into
$X$: $X \leftarrow \ldots \rightarrow Y$. These paths carry non-causal, spurious
association and have to be blocked. A set $Z$ satisfies the \emph{back-door
criterion} for $(X, Y)$ if two conditions hold: (1) no variable in $Z$ is a
descendant of $X$, and (2) $Z$ blocks every back-door path from $X$ to $Y$ given
$Z$

The key result is the adjustment formula: if $Z$ satisfies the back-door criterion,
\begin{equation}
  \Pr(Y \mid do(X)) = \sum_{z}\Pr(Y \mid X, Z = z) \Pr(Z = z) \label{eq:backdoor-adjustment}
\end{equation}

This formula is remarkable: it turns an intervention (rung 2) into a computation
over observational quantities (rung 1). Given data and the adjustment set, this
can be estimated by simple regression or matching. The back-door criterion says precisely
which variables to condition on—not "condition on everything." A common mistake
is conditioning on a descendant of $X$ (which biases the estimate) or on a
mediator (which masks part of the causal effect). The criterion is the guardrail
against these errors

% From: book_springer/lectures_source/Lesson02.2_Integrating_Causality_And_Probability_in_ML.smd:404 '* Back-Door Adjustment: A Worked Example'
\textbf{Back-Door Adjustment: A Worked Example} --- Consider estimating the causal
effect of price on sales. Advertising spend is a confounder: it can raise the price
a product can command (premium products are advertised more) \emph{and} raise
sales directly (advertising increases demand). The back-door path is $\text{Price}
\leftarrow \text{AdSpend}\to \text{Sales}$

Conditioning on advertising spend blocks this path. The adjustment formula becomes:
\begin{equation}
  \Pr(\text{Sales}\mid do(\text{Price})) = \sum_{a}\Pr(\text{Sales}\mid \text{Price}
  , \text{AdSpend}= a) \Pr(\text{AdSpend}= a) \label{eq:price-sales-adjustment}
\end{equation}

In practice, this means: estimate the effect of price on sales for each level of
advertising spend separately, then average over the distribution of advertising
spend in the population. This removes the confounding by advertising

Common mistakes in applying back-door adjustment are worth enumerating: conditioning
on a descendant of $X$ (biases the estimate by introducing collider bias),
conditioning on a mediator (masks the indirect effect), or omitting a confounder
(leaves bias in place). The back-door criterion specifies exactly which
variables satisfy these criteria, and a causal graph is essential for identifying
the right set

% Figure: Back-door DAG
% rendered_image:begin
% filename: figures/backdoor_dag.png
% label: fig:backdoor
% caption: Back-door example: Ad spend ($Z$) confounds price ($X$) and sales ($Y$). The back-door path is $X \leftarrow Z \to Y$.
% rendered_image:end

% From: book_springer/lectures_source/Lesson02.2_Integrating_Causality_And_Probability_in_ML.smd:437 '* When Back-Door Fails: The Front-Door Criterion'
\textbf{When Back-Door Fails: The Front-Door Criterion} --- Sometimes no
observable set blocks every back-door path, because the confounder itself is unobserved
In this case, the \emph{front-door criterion} offers a second path to identification,
provided the effect flows through a mediator. The front-door criterion identifies
$\Pr(Y \mid do(X))$ if:

\begin{enumerate}
  \item All paths from $X$ to $Y$ go through a mediator $M$

  \item No unobserved confounder affects both $X$ and $M$

  \item $X$ blocks every back-door path from $M$ to $Y$ (this is the technical
    condition; intuitively, the mediator is not confounded)
\end{enumerate}

If these conditions hold, the causal effect can be computed via a formula that applies
back-door adjustment twice:
\begin{equation}
  \Pr(Y \mid do(X)) = \sum_{m}\Pr(M = m \mid X) \sum_{x'}\Pr(Y \mid M = m, X = x
  ') \Pr(X = x') \label{eq:frontdoor-adjustment}
\end{equation}

Consider whether watching ads ($X$) increases cereal purchases ($Y$). An
unobserved confounding factor—how much parents care about breakfast ($U$)—could drive
both Some parents care a lot about breakfast and are more likely to buy cereal,
and they also pay attention to advertising. The back-door adjustment fails because
$U$ is unobserved

But ads work through a mediator: they make kids nag ($M$) for cereal. Parents
who care about breakfast may buy cereal regardless of ads, but ads primarily affect
purchases through nagging. If $U$ does not directly affect nagging, only ads do,
the front-door criterion holds, and we can use Equation~\eqref{eq:frontdoor-adjustment}
to identify the effect. The formula first estimates the effect of ads on nagging,
then the effect of nagging on purchases, and combines them to get the total effect
of ads on purchases. Front-door adjustment is thus a third path to
identification, available when back-door adjustment fails but a mediator is available
and appropriately separated from confounders

% Figure: Front-door DAG
% rendered_image:begin
% filename: figures/frontdoor_dag.png
% label: fig:frontdoor
% caption: Front-door example: Unobserved care about breakfast ($U$) confounds ads ($X$) and purchases ($Y$). Ads work through kids nagging ($M$), which is not affected by $U$.
% rendered_image:end

% From: book_springer/lectures_source/Lesson02.2_Integrating_Causality_And_Probability_in_ML.smd:475 '* Instrumental Variables: A Third Path'
\textbf{Instrumental Variables: A Third Path} --- An \emph{instrumental variable}
(IV) is a third mechanism for identification. An IV $Z$ is a variable that
affects the treatment $X$ but has no path to the outcome $Y$ except through $X$:

\begin{itemize}
  \item $Z$ affects $X$ (the IV has a strong causal effect on the treatment)

  \item $Z$ does not affect $Y$ directly or through any back-door path (the IV's
    only path to $Y$ is through $X$)
\end{itemize}

The idea is that $Z$ provides exogenous variation in $X$ that stands in for a randomized
trial. If $Z$ is truly random with respect to confounders, the variation in $X$ induced
by $Z$ is unconfounded

A classic example: estimating the causal effect of schooling ($X$) on earnings ($Y$)
There are likely confounders (ability, motivation) that affect both schooling
and earnings, making observational estimates biased. Distance to the nearest
college ($Z$) affects how much schooling someone obtains (cost and convenience determine
enrollment) but has no direct effect on earnings. The IV estimate uses variation
in schooling caused by distance to college, which is unconfounded by ability

Instrumental variables have limitations: a credible instrument is hard to find, and
the estimate applies only to the \emph{complier} population (units whose treatment
is actually affected by the instrument). In the schooling example, IV estimates
apply to people whose schooling decisions are affected by distance to college, not
to people who would attend college regardless of distance. Back-door adjustment,
front-door adjustment, and instrumental variables are three different ways of
turning the same problem—identifying $\Pr(Y \mid do(X))$—into a computation over
observational data. Each is suited to different graph structures and degrees of
unobserved confounding

% From: book_springer/lectures_source/Lesson02.2_Integrating_Causality_And_Probability_in_ML.smd:493 '# Potential Outcomes and Treatment Effects'
\section{Potential Outcomes and Treatment Effects}
\label{sec:potential-outcomes-treatment-effects}

% From: book_springer/lectures_source/Lesson02.2_Integrating_Causality_And_Probability_in_ML.smd:495 '## Measuring What Cannot Be Observed'
\subsection{Measuring What Cannot Be Observed}
\label{subsec:measuring-what-cannot-be-observed}

% From: book_springer/lectures_source/Lesson02.2_Integrating_Causality_And_Probability_in_ML.smd:497 '* Potential Outcomes: Rubin's Framework'
\textbf{Potential Outcomes: Rubin's Framework} --- So far, the chapter has used Pearl's
do-calculus notation. Donald Rubin's framework for causal inference uses a
different language, one that makes the problem of unobserved outcomes explicit A
unit $i$ has two \emph{potential outcomes} under a binary treatment: $Y_{1i}$ if
treated, $Y_{0i}$ if not. The relationship to Pearl's notation is $Y_{t,i}\equiv
Y_{i}\mid do(T_{i}= t)$

Only one potential outcome is ever observed, or "factual"; the other is "counterfactual."
Which one we see depends on the treatment received:
\begin{equation}
  Y_{i}= T_{i}\cdot Y_{1i}+ (1 - T_{i}) \cdot Y_{0i}\label{eq:factual-counterfactual}
\end{equation}

The \emph{fundamental problem of causal inference} is that $Y_{1i}$ and $Y_{0i}$
can never both be observed for the same unit. A business either cuts prices or does
not in a given week; the amount it would have sold under the other choice is
never recorded. An employee either receives training or does not; how they would
have performed without training is never measured

Causal inference estimates population quantities without ever observing an individual
unit's treatment effect. The average treatment effect (ATE), which we develop next,
is an example: it is the average of $Y_{1}- Y_{0}$ over the population, but we
never observe $Y_{1}- Y_{0}$ for any individual. Causal inference is about estimating
aggregate effects from data that only reveal partial information about individuals
This is possible through careful use of randomization or causal assumptions, but
it is a form of induction, not a direct measurement

% From: book_springer/lectures_source/Lesson02.2_Integrating_Causality_And_Probability_in_ML.smd:513 '* The Bias Equation'
\textbf{The Bias Equation} --- The quantity we can measure from data, $\mathbb{E}
[Y \mid T = 1] - \mathbb{E}[Y \mid T = 0]$, is not automatically the causal
effect $\mathbb{E}[Y_{1}- Y_{0}]$. The two decompose as:
\begin{equation}
  \underbrace{\mathbb{E}[Y \mid T = 1] - \mathbb{E}[Y \mid T = 0]}_{\text{observed
  difference}}= \underbrace{\mathbb{E}[Y_1 - Y_0 \mid T = 1]}_{\text{ATT}}+ \underbrace{\mathbb{E}[Y_0
  \mid T = 1] - \mathbb{E}[Y_0 \mid T = 0]}_{\text{bias}}\label{eq:bias-equation}
\end{equation}

The bias term captures whether treated and untreated units are \emph{exchangeable}
(have the same distribution of potential outcomes) but for the treatment. Association
equals causation only when the bias term vanishes

A concrete example: businesses that cut prices differ from those that do not, in
size, location, competitive pressure, and timing, before any price cut happens
If larger or more competitive firms are more likely to cut prices, and size/competition
also affect sales independently, then the observed difference in sales between
price-cutting and non-cutting firms includes both the causal effect of the cut
and the pre-existing differences. The bias term is exactly that pre-existing difference

Randomization is the one intervention guaranteed to zero out the bias term, by construction:
random assignment ensures that treated and untreated units have identical
distributions of potential outcomes (i.e., exchangeability). Without randomization,
the bias term remains unless it is explicitly addressed through stratification, matching,
or covariate adjustment. Association equals causation only under exchangeability
The bias equation shows precisely where the error lies when exchangeability is
violated, and it is the starting point for every causal inference method: adjust
for, block, or otherwise deal with the bias term

% From: book_springer/lectures_source/Lesson02.2_Integrating_Causality_And_Probability_in_ML.smd:533 '* Average and Conditional Treatment Effects'
\textbf{Average and Conditional Treatment Effects} --- Three closely related
quantities capture different aspects of the causal effect:

The \emph{average treatment effect} (ATE) is the population-level impact of treatment:
\begin{equation}
  \text{ATE}\defeq \mathbb{E}[Y_{1}- Y_{0}] = \mathbb{E}[Y \mid do(T = 1)] - \mathbb{E}
  [Y \mid do(T = 0)] \label{eq:ate}
\end{equation}

The \emph{average treatment effect on the treated} (ATT) restricts to units that
actually received treatment:
\begin{equation}
  \text{ATT}\defeq \mathbb{E}[Y_{1}- Y_{0}\mid T = 1] \label{eq:att}
\end{equation}

The \emph{conditional average treatment effect} (CATE) restricts to a subgroup defined
by covariates $X = x$:
\begin{equation}
  \text{CATE}(x) \defeq \mathbb{E}[Y_{1}- Y_{0}\mid X = x] \label{eq:cate}
\end{equation}

These quantities answer different policy questions. The ATE answers: should we roll
this intervention out to everyone? If the average effect is positive, yes. The
ATT answers: what was the effect on those who received it? (Useful for post-hoc evaluation.)
The CATE answers: which customers should get the discount? (Useful for
personalization.) In practice, CATE is the most valuable for decisions, because
it turns a single population-wide answer into a personalized policy. Estimating CATE
requires more data or stronger assumptions than ATE, but it is often worth the
extra effort because the heterogeneity in treatment effects is decision-relevant

% From: book_springer/lectures_source/Lesson02.2_Integrating_Causality_And_Probability_in_ML.smd:551 '* Assumptions Are Not Free'
\textbf{Assumptions Are Not Free} --- Every identification result in this
lesson—back-door adjustment, front-door adjustment, instrumental variables, and the
bias equation—rests on assumptions that do not verify themselves from data alone
Two assumptions are worth emphasizing

\emph{Positivity} (or common support) requires that every covariate group has a
positive chance of receiving each treatment: $0 < \Pr(T \mid X) < 1$. Without
positivity, some units have no comparable counterfactual in the data at all. If
all wealthy customers received the discount, there is no way to estimate the effect
of the discount on wealthy customers from data alone

\emph{Ignorability} (or unconfoundedness) requires that, given the adjustment
set $X$, potential outcomes are independent of treatment:
$(Y_{0}, Y_{1}) \perp T \mid X$. This assumption fails silently when a relevant confounder
is missing from the adjustment set. It is impossible to check from data whether
an unobserved confounder exists; domain knowledge and causal reasoning are required

A causal estimate is only as credible as the graph and the assumptions behind it
Correct arithmetic on a wrong graph still produces a wrong answer. Sensitivity analysis
(examining how robust conclusions are to violations of assumptions) is a crucial
complement to identification theory

% From: book_springer/lectures_source/Lesson02.2_Integrating_Causality_And_Probability_in_ML.smd:569 '# Toward Probabilistic Decision-Making'
\section{Toward Probabilistic Decision-Making}
\label{sec:toward-probabilistic-decision-making}

% From: book_springer/lectures_source/Lesson02.2_Integrating_Causality_And_Probability_in_ML.smd:571 '## Why Point Estimates Are Not Enough'
\subsection{Why Point Estimates Are Not Enough}
\label{subsec:why-point-estimates-are-not-enough}

The identification machinery developed above computes $\Pr(Y \mid do(X))$ from
data. But a point estimate—a single number for the causal effect—hides how much a
decision-maker could be wrong. This section introduces the probabilistic toolkit:
methods for attaching a credible range to a causal estimate

% From: book_springer/lectures_source/Lesson02.2_Integrating_Causality_And_Probability_in_ML.smd:573 '* Posteriors, Not Point Estimates'
\textbf{Posteriors, Not Point Estimates} --- A point estimate conflates two distinct
kinds of uncertainty. \emph{Aleatoric uncertainty} is irreducible randomness:
even with infinite data, some variation in outcomes is due to noise in the world,
not to lack of knowledge. \emph{Epistemic uncertainty} is ignorance: with more data,
this part shrinks. A point estimate, giving one number, obscures this distinction

A \emph{Bayesian} treatment replaces a point estimate $\hat{\theta}$ with a posterior
distribution:
\begin{equation}
  \Pr(\theta \mid \mathcal{D}) \propto \Pr(\mathcal{D}\mid \theta) \Pr(\theta) \label{eq:bayes-rule}
\end{equation}

The posterior's spread is the epistemic uncertainty; a separate noise term in
the structural equation captures the aleatoric part. If the posterior is very wide,
it means the data are not very informative about the parameter. If the posterior
is narrow, it means the data have pinned down the parameter tightly. A causal effect
estimated from a back-door or front-door adjustment is not one number, it is a
posterior over that number, and the width of that posterior is itself decision-relevant
A wide posterior says: "I do not have enough data to be confident"; a narrow posterior
says: "The data have resolved this question."

% From: book_springer/lectures_source/Lesson02.2_Integrating_Causality_And_Probability_in_ML.smd:587 '* Variational Inference: Scaling Bayesian Updating'
\textbf{Variational Inference: Scaling Bayesian Updating} --- Computing the exact
posterior $\Pr(\theta \mid \mathcal{D})$ in closed form is rarely possible once the
model has real complexity. \emph{Variational inference} replaces exact inference
with optimization. The idea is to pick a tractable family of distributions
$q_{\phi}(\theta)$ (e.g., a Gaussian) and fit $\phi$ to minimize the distance to
the true posterior. The objective is the evidence lower bound (ELBO):
\begin{equation}
  \text{ELBO}(\phi) = \mathbb{E}_{q_\phi}[\log \Pr(\mathcal{D}\mid \theta)] - \text{KL}
  (q_{\phi}(\theta) \| \Pr(\theta)) \label{eq:elbo}
\end{equation}

Maximizing ELBO finds the approximate posterior that best matches the true posterior
The first term rewards good fit to data; the second term regularizes the fit to stay
close to the prior. Variational inference turns "compute the posterior" into "solve
an optimization problem," which scales to the large models this book eventually needs—a
practical necessity for Bayesian inference in realistic settings

% From: book_springer/lectures_source/Lesson02.2_Integrating_Causality_And_Probability_in_ML.smd:603 '* Normalizing Flows: Flexible Posteriors'
\textbf{Normalizing Flows: Flexible Posteriors} --- A variational family that is
too simple—for example, a single Gaussian—can badly underfit a posterior with
skew or multiple modes. A \emph{normalizing flow} builds a flexible distribution
by chaining invertible transformations of a simple base distribution. Each transformation
is invertible, so exact density evaluation and sampling are possible:
\begin{equation}
  q(z) = q_{0}(z_{0}) \prod_{t=1}^{T}\left| \det \frac{\partial \phi_{t}}{\partial
  z_{t-1}}\right|^{-1}\label{eq:flow}
\end{equation}

where $z_{0}\sim q_{0}$ (e.g., standard normal) and $z = z_{T}$ is the final
sample. The jacobian determinant term accounts for the change in density under each
transformation. Normalizing flows let the shape of the posterior match the shape
of the actual uncertainty, instead of forcing it into a fixed parametric form This
flexibility comes at a computational cost, but when the posterior is multimodal
or highly skewed, it is worth the investment

% From: book_springer/lectures_source/Lesson02.2_Integrating_Causality_And_Probability_in_ML.smd:618 '* Calibration and Conformal Prediction'
\textbf{Calibration and Conformal Prediction} --- A model is \emph{calibrated}
if its stated confidence matches its actual accuracy. A poorly calibrated posterior
is worse than an honest point estimate: it looks trustworthy while being wrong. For
instance, if the model says it is 90\% confident in an interval, that interval
should contain the true value 90\% of the time, not 60\% of the time

\emph{Conformal prediction} is a model-agnostic procedure that wraps any model with
a guarantee: a 90\% conformal interval contains the true outcome at least 90\%
of the time, under minimal assumptions. The procedure is:

\begin{enumerate}
  \item Train a model on a training set

  \item On a held-out validation set, compute the model's prediction error for
    each example

  \item For a new test point, get the model's prediction. Add the $0.9 \cdot n$-th
    percentile of validation errors to create an interval

  \item This interval is guaranteed to contain the true value with probability at
    least 0.9 (up to small finite-sample effects)
\end{enumerate}

Calibration and conformal prediction let a decision-maker trust the width of an interval,
not just its center. This is essential for decision-making under uncertainty,
where the decision's quality depends on whether the uncertainty estimate is
honest

% From: book_springer/lectures_source/Lesson02.2_Integrating_Causality_And_Probability_in_ML.smd:634 '* Decision Readiness Under Uncertainty'
\textbf{Decision Readiness Under Uncertainty} --- Not every causal claim is safe
to act on, even with a posterior attached. A claim is \emph{decision-ready} when
three conditions hold together:

\begin{itemize}
  \item The identification assumptions (positivity and ignorability) are plausible
    for the population being acted on

  \item The posterior's epistemic width is narrow enough that the recommended action
    is stable across its range. (If different actions are recommended depending on
    which part of the posterior is true, the claim is not yet decision-ready.)

  \item The \emph{aleatoric floor}, the part no amount of data removes, is small
    relative to the stakes of the decision
\end{itemize}

A wide posterior is not a reason to ignore uncertainty and act anyway; it is the
signal to gather more data before acting, rather than after. This is a sharp departure
from the prediction-focused mindset, where a wide confidence interval just means
"this model is uncertain, but we will use the point estimate anyway." In
decision-making, uncertainty that is wide enough to flip the recommendation is a
call for more evidence, not a reason to proceed. Identification and uncertainty quantification
are necessary for a good decision, but not sufficient; the business objective and
the stakes still have to enter the picture

% From: book_springer/lectures_source/Lesson02.2_Integrating_Causality_And_Probability_in_ML.smd:650 '# Roadmap'
\section{Roadmap}
\label{sec:roadmap}

% From: book_springer/lectures_source/Lesson02.2_Integrating_Causality_And_Probability_in_ML.smd:652 '* Causal Models as the Foundation for Decisions'
\textbf{Causal Models as the Foundation for Decisions} --- A causal model built
with the tools in this lesson supplies exactly what a decision needs and a
predictive model cannot. It delivers the interventional distribution $\Pr(Y \mid
do(a))$, identified through one of the three paths developed above; combined with
a posterior over that distribution from the probabilistic toolkit, it gives a
decision-maker the full picture of what an action would do and how confident
they should be

The difference from prediction is stark. A marketing team can observe that
people who open a promotional email buy more. But this is rung-1 reasoning. Only
a causal model, identified through the right adjustment set or other identification
strategy, answers whether \emph{sending} the email causes the purchase. A
recommender system can predict who will click an ad; a causal model can answer whether
showing the ad will cause a click (which is often negative, because ads often target
people already predisposed to the action). For decisions, the interventional
layer is not an optional detail. It is the layer prediction alone cannot supply A
decision system without causal reasoning is a system that cannot distinguish whether
correlation reflects causation, confounding, or reverse causation—the three central
hazards of decision-making the previous chapter identified

% From: book_springer/lectures_source/Lesson02.2_Integrating_Causality_And_Probability_in_ML.smd:666 '* From Identification to Decisions'
\textbf{From Identification to Decisions} --- The pipeline built in this lesson feeds
directly into the next step of the book. The flow is:
\begin{equation}
  \text{data}\to \text{causal DAG}\to \text{identification}\to \Pr(Y|do(a)) \to \text{posterior}
  \to \arg\max_{a}\mathbb{E}[U(Y)|do(a)] \label{eq:causal-pipeline}
\end{equation}

Everything before the last arrow is this lesson; the last arrow, turning an
identified, uncertain effect into a choice of action, is decision theory. Identification
answers the question: ``Can we know the effect?'' Decision theory, to come next,
answers: ``Given what we know, what should we do?''

% From: book_springer/lectures_source/Lesson02.2_Integrating_Causality_And_Probability_in_ML.smd:681 '* Roadmap: What Comes Next'
\textbf{Roadmap: What Comes Next} --- This lesson turned the ladder of causation
into computable identification results, and turned point estimates into posteriors
The next lesson closes Part~I by turning those interventional, uncertain effects
into decisions. It will develop:

\begin{itemize}
  \item \emph{Utility functions and expected utility under $do(a)$}: How to encode
    what success looks like for a decision

  \item \emph{Cost asymmetry and multi-objective trade-offs}: Why the cost of a false
    positive might differ from the cost of a false negative, and how to navigate
    competing objectives

  \item \emph{Dynamics}: What happens when the world responds to the decision itself,
    when the action changes the distribution of future data
\end{itemize}

The key insight remains: identification and uncertainty quantification are necessary
for a good decision, but not sufficient. The business objective still has to enter
the picture. A decision system that can identify effects but has no way to
specify what success looks like will optimize for the wrong thing. The next lesson
supplies that missing piece
